\documentclass[10pt,a4paper]{amsart}
\usepackage[margin=19mm]{geometry}
\usepackage{amsmath,amssymb,mathtools}
\usepackage[scr=rsfs,cal=euler]{mathalfa}
\usepackage{xfrac}
\usepackage[unicode,hidelinks,bookmarks=false]{hyperref}
\newcommand{\ie}{i.e.\ }
\newcommand{\cf}{cf.\ }
\newcommand{\resp}{respectively\ }
\newcommand{\Subsection}{\subsection}
\providecommand{\nobreakdash}{\nobreak\hbox{-}}
\setlength{\emergencystretch}{2em}
\multlinegap0pt
\usepackage{amssymb}
\usepackage{mathtools}
\usepackage{xcolor}
\newtheorem{lemma}{Lemma}
\newcommand{\N}{\mathbb N}
\newcommand{\R}{\mathbb R}
\renewcommand{\d}{\mathrm d}
\title{Sharp regularity and small ball probabilities for the stochastic heat equation on bounded domains}
\author{Jingwu Hu, Cheuk Yin Lee}
\date{Source: arXiv:2609.08718v1 (CC BY 4.0)}
\begin{document}
\maketitle

\noindent\emph{Source and licence.} Sharp regularity and small ball probabilities for the stochastic heat equation on bounded domains. Version arXiv:2609.08718v1 (CC BY 4.0), distributed by arXiv under the Creative Commons Attribution 4.0 International licence, http://creativecommons.org/licenses/by/4.0/, as recorded in the arXiv licence metadata for this exact version and shown on the abstract page https://arxiv.org/abs/2609.08718v1. Full licence notice: \texttt{licenses/arxiv-2609.08718-CC-BY-4.0.txt}.\par\smallskip

\noindent\textit{Editorial scope.} Counted unit: the article's lemma lem:curve (source0.tex lines 1400--1405). The article's complete original proof of it is delivered with it (source0.tex lines 1482--1555, one \texttt{proof} environment of 74 lines, which the article places away from the statement). No mathematics is written, repaired or re-derived: the statement and the proof are the article's own lines, in the article's own notation, and no retained line is substituted or re-flowed.  No further source-local context is needed: every cross-reference of every retained line resolves inside the two delivered spans.  Nothing was omitted from any retained span, so the package carries no omission record.  No retained line contains a citation, so the wrapper carries no bibliography and no bibliography entry.  No retained line contains a figure environment, an image-inclusion command or drawing code, so no image is delivered and the package declares no asset dependency.  Editorial formatting only, in addition: an AMS \texttt{amsart} preamble that restates the article's own declarations and macros used by the retained lines, namely the environments \texttt{lemma} on separate counters, together with the standard packages those lines need.  The retained environments print in this document as Lemma 1 in order of appearance, whereas the article prints the same items with the numbering of its own full document.  The complete native source of the article as distributed in the arXiv e-print for this exact version is bundled unchanged as \texttt{sources/209740/source0.tex}.
\begin{lemma}\label{lem:curve}
If $D\subset \R^d$ is a bounded $C^k$ domain for some $k \in \N_+$, then there exist $\varepsilon_0 \in (0,1)$ and $C_0>0$ such that the following property holds: For every pair $x,y \in D$ with $|x-y|\le \varepsilon_0$, there is a $C^k$ curve $\gamma:[0,1] \to D$ such that $\gamma(0)=x$, $\gamma(1)=y$, and
\[
	|\gamma'(\tau)| \le C_0|x-y| \qquad \forall \tau \in [0,1].
\]
\end{lemma}

\begin{proof}[Proof of Lemma \ref{lem:curve}]
Since $D$ is a $C^k$ domain, for every $q \in \partial D$, we can find a radius $\varepsilon \in (0,1)$, a $C^k$ function $\phi:\R^{d-1} \to \R$, and an orthonormal coordinate system $(z,v)_q$ for $\R^d$, with $z \in \R^{d-1}, v \in \R$ and with origin $(0,0)_q=q$ such that
\begin{align*}
	&D \cap B(q,\varepsilon) = B(q,\varepsilon) \cap \{ (z,v+\phi(z))_q : z \in \R^{d-1}, v>0 \},\\
	&\partial D \cap B(q,\varepsilon) = B(q,\varepsilon) \cap \{ (z, \phi(z))_q : z \in \R^{d-1} \}.
\end{align*}
Since $D$ is bounded, its boundary $\partial D$ is compact, so we can find $q_1,\dots,q_n \in \partial D$ such that $\partial D \subset \bigcup_{i=1}^n B(q_i, \varepsilon_i/12)$, where $\varepsilon_1,\dots,\varepsilon_n \in (0,1)$ are the corresponding radii, $\phi_1,\dots, \phi_n :\R^{d-1} \to \R$ are the corresponding $C^1$ functions, and $(z,v)_i$ are the corresponding coordinate systems.

Since each $\phi_i$ is $C^k$ for some $k \ge 1$,
\begin{align}\label{curve:L}
	L_i := \sup_{|z| \le \varepsilon_i} |\nabla \phi_i(z)| < \infty.
\end{align}
We may take 
\begin{align}\label{curve:eps_0}
	C_0 := 2\left[1+\max_{1 \le i \le n}L_i^2\right]
	\quad \text{and}\quad
	\varepsilon_0 := \min_{1\le i\le n} \left[ \frac{\varepsilon_i}{12} \wedge \frac{\varepsilon_i}{4C_0} \right].
\end{align}
Consider any pair $x, y \in D$ with $|x-y| \le \varepsilon_0$.
Since $\overline{B}(x,\varepsilon_0) \cap \bigcup_{i=1}^n B(q_i, \varepsilon_i/12)$ is either empty or nonempty, and since $\varepsilon_0 \le \varepsilon_i/12$, either one of the following holds:
\begin{enumerate}
\item $\overline{B}(x, \varepsilon_0) \subset D$, or
\item $\overline{B}(x, \varepsilon_0) \subset B(q_i, \varepsilon_i/4)$ for some $i \in \{1,\dots, n\}$.
\end{enumerate}

In case (1), we may take $\gamma$ as the straight line from $x$ to $y$, and get $|\gamma'(\tau)| \le |x-y|$ for all $\tau \in [0,1]$.
In case (2), we may write 
\[
	x = F(z_1,v_1) \quad \text{and} \quad 
	y = F(z_2,v_2)
\]
for some $z_1, z_2 \in \R^{d-1}$ and $v_1,v_2>0$, where $F: \R^{d-1} \times \R \to \R^d$ is given by
\[
	F(z,v) = (z, v+ \phi_i(z))_i.
\]
Define $\gamma: [0,1] \to \R^d$ by
\[
	\gamma(\tau) = (\tau z_2+(1-\tau)z_1 \ ,\  \tau v_2+(1-\tau)v_1 + \phi_i(\tau z_2+(1-\tau)z_1))_i.
\]
Then, $\gamma$ is a $C^k$ curve such that $\gamma(0) = x$, $\gamma(1)=y$, and
\begin{align*}
	\gamma'(\tau) &= ( z_2-z_1\ , \ v_2-v_1 + \nabla \phi_i(\tau z_2+(1-\tau)z_1) \cdot (z_2-z_1) )_i.
\end{align*}
By \eqref{curve:L}, we have
\begin{align}\begin{split}\label{gamma':bd}
	|\gamma'(\tau)| &\le \sqrt{ |z_1-z_2|^2 + 2|v_1-v_2|^2+2L_i^2 |z_1-z_2|^2}\\
	& \le \sqrt{2(1+L_i^2)}\, |(z_1,v_1)-(z_2,v_2)|
\end{split}\end{align}
for all $\tau \in [0,1]$.
Note that $F$ is invertible with $F^{-1}(a,b) = (a,b-\phi_i(a))$ and
\begin{align*}
	|F^{-1}(a_1,b_1) - F^{-1}(a_2,b_2)|
	&\le \sqrt{2(1+L_i^2)}\, |(a_1,b_1)-(a_2,b_2)|
\end{align*}
for all $(a_1,b_1), (a_2,b_2) \in \R^{d-1} \times \R$.
It follows that
\[
	|(z_1,v_1)-(z_2,v_2)| = |F^{-1}(x)-F^{-1}(y)| \le \sqrt{2(1+L_i^2)} \, |x-y|.
\]
The preceding together with \eqref{gamma':bd} and the choice of $C_0$ in \eqref{curve:eps_0} implies that
\begin{align}\label{gamma':bd2}
	|\gamma'(\tau)| \le 2(1+L_i^2) |x-y| \le C_0|x-y| \quad \text{for all $\tau \in [0,1]$.}
\end{align}
This yields the desired estimate.
It remains to verify that the curve $\gamma$ lies in $D$.
Indeed, \eqref{gamma':bd2}, $|x-y| \le \varepsilon_0$, and $\overline{B}(x,\varepsilon_0) \subset B(q_i,\varepsilon_i/4)$ imply that for every $\tau \in [0,1]$, 
\begin{align*}
	|\gamma(\tau)-q_i| &\le |\gamma(\tau)-x| + |x-q_i| = \left|\int_0^\tau \gamma'(r) \, \d r\right| + |x-q_i|\\
	&\le C_0 \varepsilon_0 + \frac{\varepsilon_i}{4}
	\le C_0\left(\frac{\varepsilon_i}{4C_0}\right) + \frac{\varepsilon_i}{4} = \frac{\varepsilon_i}{2},
\end{align*}
where the last inequality follows from the choice of $\varepsilon_0$ in \eqref{curve:eps_0}. This shows that $\gamma(\tau) \in \overline{B}(q_i,\varepsilon_i/2) \cap \{ (z,v+\phi_i(z))_i : z \in \R^{d-1}, v>0 \} \subset D \cap \overline{B}(q_i,\varepsilon_i/2)$ for all $\tau \in [0,1]$.
The proof of Lemma \ref{lem:curve} is thus complete.
\end{proof}

\end{document}
